\documentclass[11pt,a4paper]{article}

\usepackage[spanish,es-noshorthands,es-nodecimaldot]{babel}
\usepackage[T1]{fontenc}
\usepackage{lmodern}
\usepackage{microtype}
\usepackage[margin=2.2cm]{geometry}
\usepackage{amsmath,amssymb}
\usepackage{xcolor}
\usepackage{tikz}
\usetikzlibrary{arrows.meta,positioning,shapes.geometric}
\usepackage[most]{tcolorbox}
\usepackage{booktabs,array}
\usepackage{enumitem}
\usepackage[hidelinks]{hyperref}

% ---------- Colores (blanco y negro) ----------
\definecolor{tinta}{gray}{0}
\definecolor{nota}{gray}{0}
\definecolor{entra}{gray}{0}
\definecolor{sale}{gray}{0}
\definecolor{azul}{gray}{0}
\definecolor{resalta}{gray}{0.82}
\definecolor{region}{gray}{0.80}

% Estilos de línea para distinguir restricciones sin usar color
\tikzset{
  r1/.style={thick},
  r2/.style={thick, dashed},
  r3/.style={very thick, dotted},
  r4/.style={thick, dash dot},
  zline/.style={thick, gray!55},
}
\newcommand{\swatch}[1]{\tikz[baseline=-0.6ex]\draw[#1] (0,0)--(0.6,0);}

% ---------- Atajos ----------
\newcommand{\nota}[1]{\qquad\text{\small\color{nota}\itshape$\leftarrow$ #1}}
\newcommand{\entra}[1]{\textcolor{entra}{\textbf{#1}}}
\newcommand{\sale}[1]{\textcolor{sale}{\textbf{#1}}}
\newcommand{\res}[1]{\colorbox{resalta}{$\displaystyle #1$}}
\newcommand{\Binv}{B^{-1}}
\newcommand{\paso}[1]{\par\medskip\noindent{\color{tinta}\bfseries\large #1}\par\smallskip}
\renewcommand{\arraystretch}{1.25}
\setlength{\parindent}{0pt}
\setlength{\parskip}{5pt}

\newtcolorbox{idea}[1][]{enhanced, breakable, colback=nota!6, colframe=nota!70!black,
  arc=2mm, boxrule=0.6pt, left=8pt, right=8pt, top=6pt, bottom=6pt,
  fonttitle=\bfseries, title={#1}}
\newtcolorbox{ojo}{enhanced, breakable, colback=resalta!35, colframe=resalta!80!black,
  arc=2mm, boxrule=0.5pt, left=8pt, right=8pt, top=4pt, bottom=4pt}

\title{\textbf{Método Simplex Revisado}\\[4pt]\large Explicación del método y ejercicios resueltos}
\author{Investigación de Operaciones · Primer previo}
\date{}

\begin{document}
\maketitle
\tableofcontents
\newpage

% =====================================================================
\part{¿En qué consiste el Método Simplex Revisado?}

\section{¿Qué tipo de problema resuelve?}

Muchas decisiones de una empresa consisten en \textbf{repartir recursos que son limitados} (horas de trabajo,
materia prima, dinero, espacio) para \textbf{ganar lo máximo} o \textbf{gastar lo mínimo}. Cuando esos problemas
se pueden escribir usando solo sumas y multiplicaciones por números fijos, se les llama problemas de
\textbf{programación lineal}. El método simplex es la forma más conocida de resolverlos.

Para entenderlo usaremos un ejemplo pequeño, que en la Parte II se resuelve completo:

\begin{ojo}
Una fábrica hace \textbf{puertas} y \textbf{ventanas}. Cada lote de puertas deja una ganancia de 3 (miles de
dólares) y cada lote de ventanas deja 5. Los productos pasan por tres plantas con horas limitadas:
la planta 1 tiene 4 horas (cada lote de puertas usa 1), la planta 2 tiene 12 horas (cada lote de ventanas usa 2)
y la planta 3 tiene 18 horas (las puertas usan 3 y las ventanas 2). ¿Cuánto conviene fabricar de cada uno?
\end{ojo}

Todo problema de programación lineal tiene estas partes:

\begin{center}\small
\begin{tabular}{p{3.4cm}p{6.2cm}p{4.6cm}}
\toprule
\textbf{Parte} & \textbf{Qué es} & \textbf{En el ejemplo} \\
\midrule
Variables de decisión & Lo que tenemos que decidir. & $x_1$ = lotes de puertas \newline $x_2$ = lotes de ventanas \\
Función objetivo & Lo que queremos hacer lo más grande (o lo más pequeño) posible. & Ganancia: $Z = 3x_1 + 5x_2$ \\
Restricciones & Los límites que no se pueden pasar. & $x_1 \le 4$ \newline $2x_2 \le 12$ \newline $3x_1 + 2x_2 \le 18$ \\
No negatividad & No se pueden fabricar cantidades negativas. & $x_1 \ge 0,\ x_2 \ge 0$ \\
\bottomrule
\end{tabular}
\end{center}

% =====================================================================
\section{La idea del método simplex}

Si dibujamos todas las combinaciones de puertas y ventanas que respetan los límites, obtenemos una figura
(la zona gris). Cada punto de esa zona es un plan de producción posible.

\begin{center}
\begin{tikzpicture}[scale=0.55, >=Stealth]
  \fill[region] (0,0)--(4,0)--(4,3)--(2,6)--(0,6)--cycle;
  \draw[very thin, gray!30] (0,0) grid (8,9);
  \draw[->] (0,0)--(8.5,0) node[right] {$x_1$ (puertas)};
  \draw[->] (0,0)--(0,9.5) node[above] {$x_2$ (ventanas)};
  \foreach \x in {2,4,6,8} \node[below, font=\scriptsize] at (\x,0) {\x};
  \foreach \y in {2,4,6,8} \node[left, font=\scriptsize] at (0,\y) {\y};
  \draw[r1] (4,0)--(4,9);
  \draw[r2] (0,6)--(8,6);
  \draw[r3] (0,9)--(6,0);
  \draw[->, ultra thick] (0,0)--(0,5.75);
  \draw[->, ultra thick] (0,6)--(1.75,6);
  \foreach \p/\n/\pos in {(0,0)/A/below left, (4,0)/B/above right, (4,3)/C/right, (2,6)/D/above right, (0,6)/E/above right}
    \fill \p circle (3.5pt) node[\pos, font=\small\bfseries] {\n};
\end{tikzpicture}
\end{center}

Hay un resultado muy importante: \textbf{la mejor solución siempre está en una de las esquinas} de la figura
(en el dibujo, los puntos A, B, C, D y E). El problema es que en casos reales puede haber miles o millones de
esquinas, y revisarlas todas no es práctico. El método simplex hace algo más inteligente:

\begin{enumerate}[itemsep=3pt]
  \item Empieza en una esquina fácil de encontrar: normalmente \textbf{no fabricar nada} (punto A).
  \item Revisa si moverse hacia una esquina vecina \textbf{mejora} la ganancia.
  \item Si alguna mejora, se mueve hacia la que más promete.
  \item Repite hasta llegar a una esquina donde \textbf{ninguna vecina mejora}. Esa es la solución óptima.
\end{enumerate}

En el ejemplo el método recorre A $\to$ E $\to$ D (las flechas del dibujo): empieza sin fabricar nada
($Z = 0$), luego fabrica 6 lotes de ventanas ($Z = 30$) y termina fabricando 2 lotes de puertas y 6 de ventanas
($Z = 36$). No necesita visitar B ni C.

% =====================================================================
\section{¿Qué tiene de ``revisado''?}

La versión tradicional del simplex trabaja con una \textbf{tabla} que contiene todos los números del problema.
En cada paso esa tabla se \textbf{vuelve a escribir completa}, aunque para decidir el siguiente movimiento solo se
usan unos pocos números. Con un problema pequeño esto no importa, pero con miles de variables es muchísimo trabajo.

El \textbf{Simplex Revisado} visita las mismas esquinas y llega al mismo resultado, pero organiza el trabajo de
otra manera: en lugar de reescribir toda la tabla, guarda solo una matriz pequeña, llamada $\Binv$, y cada vez
que necesita un número lo calcula usando esa matriz y los datos originales del problema, que nunca se modifican.

\begin{center}\small
\begin{tabular}{p{3.6cm}p{5.2cm}p{5.4cm}}
\toprule
 & \textbf{Simplex tradicional (tabla)} & \textbf{Simplex Revisado} \\
\midrule
¿Qué se guarda? & La tabla completa. & Solo la matriz $\Binv$ y los datos originales. \\
¿Qué se actualiza en cada paso? & Todos los números de la tabla. & Solo $\Binv$. \\
¿Cuándo conviene? & Problemas pequeños, a mano. & Problemas grandes y en computadora. \\
\bottomrule
\end{tabular}
\end{center}

Una forma de imaginarlo: en vez de copiar todo el cuaderno cada vez que se avanza, se lleva una sola hoja
de resumen ($\Binv$) y, cuando hace falta un dato, se consulta el enunciado original.

% =====================================================================
\section{Antes de empezar: preparar el problema}

El método necesita que todas las restricciones sean \textbf{igualdades} ($=$). Para lograrlo se agregan
variables auxiliares:

\begin{description}[leftmargin=0pt, itemsep=5pt]
  \item[Restricción ``$\le$'': se suma una holgura.] La holgura representa \textbf{lo que sobra} del recurso.
        Por ejemplo, $x_1 \le 4$ se escribe $x_1 + s_1 = 4$: si se fabrica 1 lote de puertas, sobran $s_1 = 3$ horas.
  \item[Restricción ``$\ge$'': se resta un exceso.] El exceso representa \textbf{cuánto se pasa} del mínimo.
        Por ejemplo, $x_1 + x_2 \ge 5$ se escribe $x_1 + x_2 - e_1 = 5$.
  \item[Restricciones ``$\ge$'' y ``$=$'': se agrega además una variable artificial.] Es una variable
        \textbf{ayudante} que solo sirve para poder arrancar el método. Al final \textbf{debe valer cero}; si no
        se puede lograr, el problema no tiene solución. Hay dos formas de hacerla desaparecer: el método de la
        \textbf{Gran M} (se le pone un castigo enorme en la función objetivo) y el método de las
        \textbf{Dos Fases} (primero se eliminan las ayudantes y después se busca la mejor ganancia).
  \item[Si el problema es de minimizar,] se cambian los signos de la función objetivo y se maximiza. Al final
        se le vuelve a cambiar el signo al resultado.
\end{description}

Las holguras (y las artificiales, si las hay) son el punto de partida: al principio \textbf{no se fabrica nada
y todo el recurso sobra}. En el ejemplo el problema preparado queda:
\[
\begin{aligned}
\text{Max } Z &= 3x_1 + 5x_2 \\
x_1 + s_1 &= 4 \\
2x_2 + s_2 &= 12 \\
3x_1 + 2x_2 + s_3 &= 18
\end{aligned}
\]

% =====================================================================
\section{Las palabras que se usan}

\begin{center}\small
\begin{tabular}{p{3.2cm}p{11cm}}
\toprule
\textbf{Término} & \textbf{Qué significa} \\
\midrule
Iteración & Cada vuelta del método; en cada una se pasa de una esquina a otra. \\
Variables básicas & Las variables que están ``activas'' en la esquina actual y pueden valer más que cero.
                    Siempre hay tantas como restricciones. Al inicio son las holguras. \\
Variables no básicas & Las que valen cero en la esquina actual. \\
Base & El grupo de variables básicas. \\
$B$ y $\Binv$ & $B$ es la matriz formada por las columnas de las variables básicas; $\Binv$ es su inversa.
               Al inicio $B$ es la matriz identidad, así que $\Binv$ también. \\
$b$ & Los números del lado derecho de las restricciones (lo disponible de cada recurso). \\
$x_B$ & Cuánto vale cada variable básica. \\
$c_B$ & La ganancia de cada variable básica (las holguras tienen ganancia 0). \\
$\pi$ & Cuánto vale una unidad extra de cada recurso. También se llama \textbf{precio sombra}. \\
$z_j - c_j$ & Indica si conviene activar una variable: si es negativo, activarla aumenta la ganancia. \\
$\alpha$ & Cuánto cambia cada variable básica cuando la nueva variable aumenta en una unidad. \\
\bottomrule
\end{tabular}
\end{center}

% =====================================================================
\section{Los pasos del Simplex Revisado}\label{sec:algoritmo}

Cada iteración responde cinco preguntas, siempre en el mismo orden. Para que se entiendan, se muestra lo que
ocurre en la primera iteración del ejemplo.

\paso{Paso 1. ¿Cuánto vale cada variable básica ahora?}
Se multiplica la matriz $\Binv$ por el lado derecho $b$, y la ganancia se obtiene multiplicando la ganancia de
cada básica por su valor:
\[
x_B = \Binv\, b, \qquad Z = c_B \cdot x_B .
\]
\textit{En el ejemplo:} al inicio $\Binv$ es la identidad, así que $x_B = b$: sobran $s_1 = 4$, $s_2 = 12$ y
$s_3 = 18$ horas, y la ganancia es $Z = 0$ porque todavía no se fabrica nada.

\paso{Paso 2. ¿Cuánto vale cada recurso?}
Se multiplica la fila de ganancias de las básicas por $\Binv$:
\[
\pi = c_B \cdot \Binv .
\]
\textit{En el ejemplo:} las holguras no dan ganancia, así que $\pi = (0,\ 0,\ 0)$: por ahora ninguna hora extra
aumentaría la ganancia.

\paso{Paso 3. ¿Conviene activar alguna variable?}
Para cada variable que vale cero (no básica) se calcula
\[
z_j - c_j = \pi \cdot a_j - c_j ,
\]
donde $a_j$ es la columna de números de esa variable en las restricciones y $c_j$ su ganancia.
\begin{itemize}[itemsep=2pt]
  \item Si \textbf{ninguno es negativo}, ya no se puede mejorar: \textbf{se llegó a la solución óptima}.
  \item Si hay negativos, \entra{entra} a la base la variable con el valor \textbf{más negativo}, porque es la que
        más aumenta la ganancia por cada unidad.
\end{itemize}
\textit{En el ejemplo:} para las puertas $0 - 3 = -3$ y para las ventanas $0 - 5 = -5$. El más negativo es
$-5$, así que entran las ventanas ($x_2$).

\paso{Paso 4. ¿Qué variable sale?}
Al aumentar la variable que entra, las demás van cambiando y alguna llega a cero primero. Para saber cuál,
se calcula $\alpha = \Binv \cdot a_{\text{entra}}$ y se divide cada valor de $x_B$ entre su $\alpha$
(\textbf{solo cuando $\alpha$ es positivo}). La división más pequeña indica la variable que \sale{sale}.
\begin{itemize}[itemsep=2pt]
  \item La división más pequeña corresponde al recurso que \textbf{se agota primero}.
  \item Si ningún $\alpha$ es positivo, la variable puede crecer sin límite: el problema es \textbf{no acotado}.
\end{itemize}
\textit{En el ejemplo:} $\alpha = (0,\ 2,\ 2)$. Para $s_1$ no se divide ($\alpha = 0$); para $s_2$: $12/2 = 6$;
para $s_3$: $18/2 = 9$. La menor es 6, así que sale $s_2$: con 6 lotes de ventanas la planta 2 se queda sin horas.

\paso{Paso 5. Actualizar $\Binv$}
Como cambió una variable de la base, hay que actualizar la hoja de resumen. La forma más sencilla de hacerlo
a mano es trabajar por filas sobre $\Binv$:
\begin{enumerate}[itemsep=2pt]
  \item La fila de la variable que sale se \textbf{divide} entre su $\alpha$.
  \item A cada una de las otras filas se le \textbf{resta} su $\alpha$ multiplicado por la fila del punto anterior.
\end{enumerate}
Esto es lo mismo que calcular $\Binv_{\text{nueva}} = E \cdot \Binv$, donde $E$ es la matriz identidad con una
columna cambiada (en la Parte II se muestra en cada ejercicio). Después se vuelve al paso 1.

\textit{En el ejemplo:} la fila 2 se divide entre 2 y a la fila 3 se le resta 2 veces esa fila. Así se obtiene
la nueva $\Binv$ y comienza la iteración 1.

\begin{idea}[Resumen de los pasos]
\begin{enumerate}[leftmargin=*, itemsep=3pt]
  \item[\textbf{0.}] Preparar el problema (holguras, excesos y artificiales). Al inicio $\Binv = I$.
  \item \textbf{Valor de las básicas:} \quad $x_B = \Binv b$ \quad y \quad $Z = c_B\, x_B$.
  \item \textbf{Valor de los recursos:} \quad $\pi = c_B\, \Binv$.
  \item \textbf{¿Conviene activar alguna variable?} \quad $z_j - c_j = \pi\, a_j - c_j$.
        Si ninguno es negativo, es óptimo; si no, entra el más negativo.
  \item \textbf{¿Quién sale?} \quad $\alpha = \Binv a_{\text{entra}}$; la menor división $x_B / \alpha$ (con $\alpha > 0$).
  \item \textbf{Actualizar} $\Binv$ y volver al paso 1.
\end{enumerate}
\end{idea}

\begin{center}
\begin{tikzpicture}[node distance=6mm and 9mm, >=Stealth,
  caja/.style={draw, rounded corners=2pt, fill=gray!8, align=center, minimum width=5.6cm, inner sep=5pt, font=\small},
  decision/.style={draw, diamond, aspect=2.4, fill=gray!18, align=center, inner sep=1pt, font=\small},
  fin/.style={draw, rounded corners=8pt, fill=gray!25, align=center, minimum width=2.8cm, inner sep=6pt, font=\small}]
  \node[caja] (ini) {Preparar el problema\\ (se empieza sin fabricar nada)};
  \node[caja, below=of ini] (val) {1. Valor de las básicas y de los recursos};
  \node[decision, below=of val] (opt) {¿Conviene activar\\ alguna variable?};
  \node[fin, right=of opt] (optimo) {\textbf{Solución óptima}};
  \node[caja, below=of opt] (entra) {Entra la de valor más negativo};
  \node[decision, below=of entra] (acot) {¿Hay algún\\ $\alpha$ positivo?};
  \node[fin, right=of acot] (noacot) {\textbf{No acotado}};
  \node[caja, below=of acot] (sale) {Sale la de menor división\\ y se actualiza $\Binv$};
  \draw[->] (ini) -- (val);
  \draw[->] (val) -- (opt);
  \draw[->] (opt) -- node[above, font=\scriptsize] {no} (optimo);
  \draw[->] (opt) -- node[right, font=\scriptsize] {sí} (entra);
  \draw[->] (entra) -- (acot);
  \draw[->] (acot) -- node[above, font=\scriptsize] {no} (noacot);
  \draw[->] (acot) -- node[right, font=\scriptsize] {sí} (sale);
  \draw[->] (sale.west) -- ++(-0.9,0) |- node[pos=0.25, left, font=\scriptsize, align=center] {siguiente\\iteración} (val.west);
\end{tikzpicture}
\end{center}

% =====================================================================
\section{¿Cómo se sabe que terminó?}

\begin{description}[leftmargin=0pt, itemsep=5pt]
  \item[Solución óptima.] En el paso 3 ningún $z_j - c_j$ es negativo: ninguna esquina vecina mejora.
  \item[Problema no acotado.] En el paso 4 ningún $\alpha$ es positivo: la ganancia podría crecer sin fin.
        Normalmente significa que falta alguna restricción en el modelo.
  \item[Problema sin solución.] Al terminar, alguna variable artificial (ayudante) sigue valiendo más que cero:
        no existe ningún plan que cumpla todas las restricciones a la vez.
  \item[Varias soluciones óptimas.] Al terminar, alguna variable no básica tiene $z_j - c_j = 0$: se puede
        activar sin que la ganancia cambie, así que hay más de un plan igual de bueno.
  \item[Empates.] Si la división más pequeña da cero, el método cambia de base sin mejorar la ganancia. En casos
        muy raros podría repetir pasos sin avanzar; para evitarlo se elige siempre la variable con el número más
        bajo (regla de Bland).
\end{description}

% =====================================================================
\section{¿Qué información extra se obtiene?}

Al terminar, los valores de $\pi$ dicen \textbf{cuánto aumentaría la ganancia si se tuviera una unidad más de
cada recurso}. En el ejemplo, al final $\pi = (0,\ \tfrac32,\ 1)$:

\begin{itemize}[itemsep=2pt]
  \item Una hora más en la planta 1 no sirve de nada (ya le sobran horas): vale 0.
  \item Una hora más en la planta 2 aumentaría la ganancia en 1.5 (es decir, \$1\,500).
  \item Una hora más en la planta 3 aumentaría la ganancia en 1 (\$1\,000).
\end{itemize}

Esto ayuda a decidir, por ejemplo, en qué planta conviene invertir en horas extra. Estos mismos valores son
la solución del llamado \textbf{problema dual}.

% =====================================================================
\section{Ventajas y desventajas}

\begin{center}\small
\begin{tabular}{p{7.2cm}p{7.2cm}}
\toprule
\textbf{Ventajas} & \textbf{Desventajas} \\
\midrule
\begin{itemize}[leftmargin=*, itemsep=2pt, topsep=2pt]
  \item Usa menos memoria, porque solo guarda $\Binv$.
  \item Hace menos cálculos cuando el problema tiene muchas variables.
  \item Los datos originales nunca se modifican, así que es más difícil equivocarse al arrastrar números.
  \item Da directamente el valor de cada recurso ($\pi$).
\end{itemize}
&
\begin{itemize}[leftmargin=*, itemsep=2pt, topsep=2pt]
  \item A mano hay que multiplicar matrices en cada paso.
  \item En problemas pequeños no ahorra trabajo frente a la tabla tradicional.
  \item Hay que volver a calcular $z_j - c_j$ en cada iteración.
\end{itemize}
\\
\bottomrule
\end{tabular}
\end{center}

% =====================================================================
\newpage
\part{Ejercicios resueltos}
\setcounter{section}{0}

Los tres ejercicios siguientes aplican, uno por uno, los pasos de la sección~\ref{sec:algoritmo}. En cada iteración se
indica de dónde sale cada valor y, al final, se muestra la interpretación gráfica.

% =====================================================================
\section{Wyndor Glass Co.: puertas y ventanas}
\textit{Fuente: Hillier y Lieberman, Introducción a la Investigación de Operaciones, cap.~3.}

\subsection*{Planteamiento del problema}
Una fábrica de vidrio quiere lanzar dos productos nuevos: una \textbf{puerta de vidrio con marco de aluminio}
y una \textbf{ventana con marco de madera}. Se producen por lotes y pasan por tres plantas:

\begin{center}\small
\begin{tabular}{lccc}
\toprule
 & \textbf{Puertas} (por lote) & \textbf{Ventanas} (por lote) & \textbf{Horas por semana} \\
\midrule
Planta 1 (marcos de aluminio) & 1 h & 0 h & 4 \\
Planta 2 (marcos de madera)   & 0 h & 2 h & 12 \\
Planta 3 (vidrio y ensamble)  & 3 h & 2 h & 18 \\
\midrule
Ganancia por lote & \$3\,000 & \$5\,000 & \\
\bottomrule
\end{tabular}
\end{center}

La pregunta: \textbf{¿cuántos lotes de cada producto conviene fabricar por semana para ganar lo máximo?}

\subsection*{Formulación del modelo}
\begin{itemize}[itemsep=2pt]
  \item $x_1$ = lotes de puertas por semana, \quad $x_2$ = lotes de ventanas por semana.
  \item \textbf{Ganancia} (en miles de dólares): cada lote de puertas deja 3 y cada lote de ventanas deja 5, así que
        $Z = 3x_1 + 5x_2$.
  \item \textbf{Planta 1:} solo las puertas pasan por ahí (1 h por lote) y hay 4 h $\Rightarrow x_1 \le 4$.
  \item \textbf{Planta 2:} solo las ventanas (2 h por lote) y hay 12 h $\Rightarrow 2x_2 \le 12$.
  \item \textbf{Planta 3:} las dos pasan (3 h y 2 h) y hay 18 h $\Rightarrow 3x_1 + 2x_2 \le 18$.
  \item No se fabrican lotes negativos: $x_1, x_2 \ge 0$.
\end{itemize}

\[
\begin{aligned}
\text{Max } Z &= 3x_1 + 5x_2\\
\text{s.a. } \quad x_1 &\le 4\\
2x_2 &\le 12\\
3x_1 + 2x_2 &\le 18\\
x_1, x_2 &\ge 0
\end{aligned}
\]

\subsection*{Forma estándar}
Cada ``$\le$'' se vuelve igualdad sumándole una \textbf{holgura}: son las \textbf{horas que sobran} en cada planta.
\begin{align*}
x_1 + s_1 &= 4 \nota{$s_1$ = horas libres de la planta 1}\\
2x_2 + s_2 &= 12 \nota{$s_2$ = horas libres de la planta 2}\\
3x_1 + 2x_2 + s_3 &= 18 \nota{$s_3$ = horas libres de la planta 3}
\end{align*}
Con el orden de columnas $(x_1, x_2, s_1, s_2, s_3)$:
\[
A = \begin{bmatrix} 1 & 0 & 1 & 0 & 0\\ 0 & 2 & 0 & 1 & 0\\ 3 & 2 & 0 & 0 & 1 \end{bmatrix},\qquad
b = \begin{bmatrix} 4\\ 12\\ 18 \end{bmatrix},\qquad
c = (3,\ 5,\ 0,\ 0,\ 0).
\]
Arrancamos \textbf{sin fabricar nada}: todas las horas sobran, así que la base inicial es $\{s_1, s_2, s_3\}$.
Sus columnas en $A$ son justo la identidad.

% ---------------------------------------------------------------------
\subsection*{Iteración 0}
\paso{1. Cuánto vale cada básica}
\[
\Binv = I, \qquad x_B = \Binv b = \begin{bmatrix} 4\\ 12\\ 18\end{bmatrix}
\nota{no fabricamos nada: todas las horas están libres}
\]
\[
c_B = (0,\ 0,\ 0) \nota{las holguras no dan ganancia}, \qquad Z = 0.
\]

\paso{2. Precios duales}
\[ \pi = c_B\,\Binv = (0,\ 0,\ 0). \]

\paso{3. ¿Alguien mejora $Z$?}
Las no básicas son $x_1$ y $x_2$. Como $\pi$ es todo cero, $\pi\,a_j = 0$:
\begin{align*}
z_1 - c_1 &= 0 - 3 = -3\\
z_2 - c_2 &= 0 - 5 = \mathbf{-5} \nota{el más negativo}
\end{align*}
\entra{Entra $x_2$}: cada lote de ventanas deja 5, más que las puertas, así que es lo que más sube la ganancia.

\paso{4. ¿Quién sale?}
\[
\alpha = \Binv a_2 = \begin{bmatrix} 0\\ 2\\ 2\end{bmatrix}
\nota{horas que gasta cada lote de ventanas en las plantas 1, 2 y 3}
\]
\begin{align*}
s_1 &:\ \alpha = 0 \ \text{no cuenta} \nota{la ventana no usa la planta 1}\\
s_2 &:\ 12 / 2 = \mathbf{6} \nota{la planta 2 alcanza para 6 lotes}\\
s_3 &:\ 18 / 2 = 9 \nota{la planta 3 alcanza para 9 lotes}
\end{align*}
La menor es 6, así que \sale{sale $s_2$}: la planta 2 es la primera que se queda sin horas.

\paso{5. Actualizo $\Binv$}
Sale la fila 2 y $\alpha_2 = 2$, entonces
$\eta = \left(-\tfrac{0}{2},\ \tfrac{1}{2},\ -\tfrac{2}{2}\right) = \left(0,\ \tfrac12,\ -1\right)$:
\[
E = \begin{bmatrix} 1 & 0 & 0\\ 0 & \frac12 & 0\\ 0 & -1 & 1\end{bmatrix},
\qquad
\Binv_{\text{nueva}} = E \cdot I = \begin{bmatrix} 1 & 0 & 0\\ 0 & \frac12 & 0\\ 0 & -1 & 1\end{bmatrix}.
\]

% ---------------------------------------------------------------------
\subsection*{Iteración 1 \quad (base $\{s_1, x_2, s_3\}$)}
\paso{1. Cuánto vale cada básica}
\begin{align*}
s_1 &= 1\cdot 4 = 4 \nota{la planta 1 sigue libre}\\
x_2 &= \tfrac12 \cdot 12 = 6 \nota{ya hacemos 6 lotes de ventanas}\\
s_3 &= -1\cdot 12 + 1\cdot 18 = 6 \nota{a la planta 3 le sobran 6 horas}
\end{align*}
Con $c_B = (0,\ 5,\ 0)$: \quad $Z = 5\cdot 6 = 30$.

\paso{2. Precios duales}
\[
\pi = c_B\,\Binv = \left(0,\ 5\cdot\tfrac12,\ 0\right) = \left(0,\ \tfrac52,\ 0\right)
\nota{por ahora, una hora más en la planta 2 valdría 2.5 mil}
\]

\paso{3. ¿Alguien mejora $Z$?}
Las no básicas ahora son $x_1$ y $s_2$:
\begin{align*}
z_1 - c_1 &= \pi\cdot(1,\,0,\,3) - 3 = 0 + 0 + 0 - 3 = \mathbf{-3}\\
z_{s_2} - c_{s_2} &= \pi\cdot(0,\,1,\,0) - 0 = \tfrac52 \nota{positivo: no conviene}
\end{align*}
\entra{Entra $x_1$}: fabricar puertas todavía sube la ganancia.

\paso{4. ¿Quién sale?}
\[
\alpha = \Binv a_1 = \begin{bmatrix} 1 & 0 & 0\\ 0 & \frac12 & 0\\ 0 & -1 & 1\end{bmatrix}
\begin{bmatrix} 1\\ 0\\ 3\end{bmatrix} = \begin{bmatrix} 1\\ 0\\ 3\end{bmatrix}
\]
\begin{align*}
s_1 &:\ 4/1 = 4\\
x_2 &:\ \alpha = 0 \ \text{no cuenta}\\
s_3 &:\ 6/3 = \mathbf{2} \nota{con 2 lotes de puertas se acaban las 6 h libres}
\end{align*}
\sale{Sale $s_3$}: la planta 3 se llena.

\paso{5. Actualizo $\Binv$}
Sale la fila 3 con $\alpha_3 = 3$: \ $\eta = \left(-\tfrac13,\ 0,\ \tfrac13\right)$.
\[
E = \begin{bmatrix} 1 & 0 & -\frac13\\ 0 & 1 & 0\\ 0 & 0 & \frac13\end{bmatrix},
\qquad
\Binv_{\text{nueva}} = E\,\Binv = \begin{bmatrix} 1 & \frac13 & -\frac13\\ 0 & \frac12 & 0\\ 0 & -\frac13 & \frac13\end{bmatrix}.
\]
Por filas: la fila 3 $(0,-1,1)$ se divide entre 3 y queda $(0,-\tfrac13,\tfrac13)$; a la fila 1 le resto
1 vez esa fila y queda $(1,\tfrac13,-\tfrac13)$; la fila 2 tiene $\alpha=0$, así que no cambia.

% ---------------------------------------------------------------------
\subsection*{Iteración 2 \quad (base $\{s_1, x_2, x_1\}$)}
\paso{1. Cuánto vale cada básica}
\begin{align*}
s_1 &= 4 + \tfrac13(12) - \tfrac13(18) = 4 + 4 - 6 = 2\\
x_2 &= \tfrac12(12) = 6\\
x_1 &= -\tfrac13(12) + \tfrac13(18) = -4 + 6 = 2
\end{align*}
Con $c_B = (0,\ 5,\ 3)$: \quad $Z = 0\cdot 2 + 5\cdot 6 + 3\cdot 2 = 36$.

\paso{2. Precios duales}
\begin{align*}
\pi_1 &= 0\cdot 1 + 5\cdot 0 + 3\cdot 0 = 0\\
\pi_2 &= 0\cdot\tfrac13 + 5\cdot\tfrac12 + 3\cdot\left(-\tfrac13\right) = \tfrac52 - 1 = \tfrac32\\
\pi_3 &= 0\cdot\left(-\tfrac13\right) + 5\cdot 0 + 3\cdot\tfrac13 = 1
\end{align*}

\paso{3. ¿Alguien mejora $Z$?}
Las no básicas son $s_2$ y $s_3$ (sus columnas son $e_2$ y $e_3$, así que $\pi\,a_j$ es simplemente $\pi_2$ y $\pi_3$):
\[
z_{s_2} - 0 = \tfrac32 \ge 0, \qquad z_{s_3} - 0 = 1 \ge 0.
\]
\textbf{Ninguno es negativo $\Rightarrow$ llegamos al óptimo.}

\subsection*{Resultado}
\[
\res{x_1 = 2,\quad x_2 = 6,\quad Z^* = 36}
\]
\begin{itemize}[itemsep=2pt]
  \item Conviene fabricar \textbf{2 lotes de puertas y 6 lotes de ventanas} por semana, ganando \textbf{\$36\,000}.
  \item $s_1 = 2$: a la planta 1 le sobran 2 horas. Las plantas 2 y 3 quedan al tope ($s_2 = s_3 = 0$).
  \item \textbf{Precios sombra} $\pi = (0,\ \tfrac32,\ 1)$: una hora extra en la planta 1 no sirve de nada
        (ya sobran horas); en la planta 2 valdría \$1\,500 y en la planta 3, \$1\,000.
\end{itemize}

\textbf{Comprobación:} $2 \le 4$ \checkmark, \ $2(6) = 12 \le 12$ \checkmark, \ $3(2) + 2(6) = 18 \le 18$ \checkmark.

\subsection*{Interpretación gráfica}
\begin{center}
\begin{tikzpicture}[scale=0.6, >=Stealth]
  \fill[region!70] (0,0)--(4,0)--(4,3)--(2,6)--(0,6)--cycle;
  \draw[very thin, gray!25] (0,0) grid (9,10);
  \draw[->] (0,0)--(9.5,0) node[right] {$x_1$};
  \draw[->] (0,0)--(0,10.5) node[above] {$x_2$};
  \foreach \x in {2,4,6,8} \node[below, font=\scriptsize] at (\x,0) {\x};
  \foreach \y in {2,4,6,8,10} \node[left, font=\scriptsize] at (0,\y) {\y};
  \draw[r1] (4,0)--(4,10);
  \draw[r2] (0,6)--(9,6);
  \draw[r3] (0,9)--(6,0);
  \draw[zline] (0,7.2)--(9,1.8);
  \draw[->, ultra thick, tinta!80] (0,0)--(0,5.8);
  \draw[->, ultra thick, tinta!80] (0,6)--(1.8,6);
  \foreach \p in {(0,0),(4,0),(4,3),(0,6)} \fill[tinta] \p circle (3pt);
  \fill[resalta!90!black] (2,6) circle (6pt);
  \node[anchor=south west, fill=white, inner sep=1pt, font=\scriptsize\bfseries, xshift=3pt, yshift=3pt] at (0,0) {It.~0};
  \node[anchor=south west, fill=white, inner sep=1pt, font=\scriptsize\bfseries, xshift=3pt, yshift=3pt] at (0,6) {It.~1};
  \node[anchor=south west, fill=white, inner sep=1pt, font=\small\bfseries, xshift=5pt, yshift=5pt] at (2,6) {It.~2: $(2,\,6)$ óptimo};
  \node[anchor=north east, draw=gray!50, fill=white, align=left, font=\scriptsize, inner sep=4pt] at (9,9.9) {
    \swatch{r1} R1: $x_1 \le 4$\\
    \swatch{r2} R2: $2x_2 \le 12$\\
    \swatch{r3} R3: $3x_1 + 2x_2 \le 18$\\
    \swatch{zline} $Z = 36$};
\end{tikzpicture}
\end{center}
El simplex arrancó en $(0,0)$, subió a $(0,6)$ al meter ventanas y terminó en $(2,6)$ al meter puertas.
Si evaluamos $Z$ en todas las esquinas: $(0,0)\to 0$, \ $(4,0)\to 12$, \ $(4,3)\to 27$, \ $(0,6)\to 30$, \
$(2,6)\to \mathbf{36}$. Coincide.

% =====================================================================
\newpage
\section{Reddy Mikks: pinturas}
\textit{Fuente: Taha, Investigación de Operaciones, cap.~2.}

\subsection*{Planteamiento del problema}
Reddy Mikks fabrica \textbf{pintura para exteriores} y \textbf{pintura para interiores} con dos materias primas,
M1 y M2:

\begin{center}
\begin{tabular}{lccc}
\toprule
 & \textbf{Exteriores} (por ton) & \textbf{Interiores} (por ton) & \textbf{Disponible por día} \\
\midrule
Materia prima M1 & 6 ton & 4 ton & 24 ton \\
Materia prima M2 & 1 ton & 2 ton & 6 ton \\
\midrule
Ganancia por tonelada & \$5\,000 & \$4\,000 & \\
\bottomrule
\end{tabular}
\end{center}

Además, el estudio de mercado dice dos cosas:
\begin{itemize}[itemsep=2pt]
  \item La producción diaria de pintura de interiores \textbf{no puede superar} a la de exteriores en más de 1 tonelada.
  \item La demanda máxima de pintura de interiores es de \textbf{2 toneladas por día}.
\end{itemize}

\subsection*{Formulación del modelo}
\begin{itemize}[itemsep=2pt]
  \item $x_1$ = toneladas diarias de pintura de exteriores, \quad $x_2$ = toneladas diarias de interiores.
  \item \textbf{Ganancia} (miles de dólares): $Z = 5x_1 + 4x_2$.
  \item \textbf{M1:} 6 ton por cada ton de exteriores y 4 por cada ton de interiores, con 24 disponibles
        $\Rightarrow 6x_1 + 4x_2 \le 24$.
  \item \textbf{M2:} $x_1 + 2x_2 \le 6$.
  \item \textbf{Mercado 1:} ``interiores no supera a exteriores en más de 1'' se escribe $x_2 \le x_1 + 1$.
        Paso $x_1$ al lado izquierdo y queda $-x_1 + x_2 \le 1$.
  \item \textbf{Mercado 2:} $x_2 \le 2$.
\end{itemize}

\[
\begin{aligned}
\text{Max } Z &= 5x_1 + 4x_2\\
\text{s.a. } \quad 6x_1 + 4x_2 &\le 24\\
x_1 + 2x_2 &\le 6\\
-x_1 + x_2 &\le 1\\
x_2 &\le 2\\
x_1, x_2 &\ge 0
\end{aligned}
\]

\subsection*{Forma estándar}
\begin{align*}
6x_1 + 4x_2 + s_1 &= 24 \nota{$s_1$ = toneladas de M1 que sobran}\\
x_1 + 2x_2 + s_2 &= 6 \nota{$s_2$ = toneladas de M2 que sobran}\\
-x_1 + x_2 + s_3 &= 1 \nota{$s_3$ = ``margen'' que queda en el mercado 1}\\
x_2 + s_4 &= 2 \nota{$s_4$ = demanda de interiores sin cubrir}
\end{align*}
\[
b = (24,\ 6,\ 1,\ 2), \qquad c = (5,\ 4,\ 0,\ 0,\ 0,\ 0), \qquad \text{base inicial } \{s_1, s_2, s_3, s_4\},\ \Binv = I.
\]

% ---------------------------------------------------------------------
\subsection*{Iteración 0}
\paso{1 y 2. Valores y precios duales}
\[
x_B = (24,\ 6,\ 1,\ 2) \nota{no producimos nada: todo sobra}, \qquad \pi = (0,0,0,0), \qquad Z = 0.
\]

\paso{3. ¿Alguien mejora $Z$?}
\[
z_1 - c_1 = 0 - 5 = \mathbf{-5}, \qquad z_2 - c_2 = 0 - 4 = -4.
\]
\entra{Entra $x_1$} (la pintura de exteriores deja más ganancia).

\paso{4. ¿Quién sale?}
\[
\alpha = \Binv a_1 = (6,\ 1,\ -1,\ 0)
\]
\begin{align*}
s_1 &:\ 24/6 = \mathbf{4} \nota{M1 alcanza para 4 ton de exteriores}\\
s_2 &:\ 6/1 = 6\\
s_3 &:\ \alpha = -1 \ \text{no cuenta} \nota{producir más exteriores \emph{afloja} esta restricción}\\
s_4 &:\ \alpha = 0 \ \text{no cuenta}
\end{align*}
\sale{Sale $s_1$}: la materia prima M1 es la que se acaba primero.

\paso{5. Actualizo $\Binv$}
Sale la fila 1 con $\alpha_1 = 6$: \
$\eta = \left(\tfrac16,\ -\tfrac16,\ -\tfrac{-1}{6},\ -\tfrac06\right) = \left(\tfrac16,\ -\tfrac16,\ \tfrac16,\ 0\right)$.
\[
\Binv_{\text{nueva}} = E = \begin{bmatrix}
\frac16 & 0 & 0 & 0\\ -\frac16 & 1 & 0 & 0\\ \frac16 & 0 & 1 & 0\\ 0 & 0 & 0 & 1
\end{bmatrix}
\]

% ---------------------------------------------------------------------
\subsection*{Iteración 1 \quad (base $\{x_1, s_2, s_3, s_4\}$)}
\paso{1. Cuánto vale cada básica}
\begin{align*}
x_1 &= \tfrac16(24) = 4\\
s_2 &= -\tfrac16(24) + 6 = -4 + 6 = 2 \nota{sobran 2 ton de M2}\\
s_3 &= \tfrac16(24) + 1 = 5\\
s_4 &= 2
\end{align*}
Con $c_B = (5,0,0,0)$: \quad $Z = 5\cdot 4 = 20$.

\paso{2. Precios duales}
\[
\pi = c_B\,\Binv = \left(5\cdot\tfrac16,\ 0,\ 0,\ 0\right) = \left(\tfrac56,\ 0,\ 0,\ 0\right).
\]

\paso{3. ¿Alguien mejora $Z$?}
La columna de $x_2$ es $a_2 = (4,\ 2,\ 1,\ 1)$:
\begin{align*}
z_2 - c_2 &= \tfrac56(4) + 0(2) + 0(1) + 0(1) - 4 = \tfrac{10}{3} - 4 = \mathbf{-\tfrac23}\\
z_{s_1} - c_{s_1} &= \pi_1 - 0 = \tfrac56 \nota{positivo: no conviene devolver M1}
\end{align*}
\entra{Entra $x_2$}: todavía conviene producir pintura de interiores.

\paso{4. ¿Quién sale?}
\[
\alpha = \Binv a_2 =
\begin{bmatrix} \frac16(4)\\[2pt] -\frac16(4) + 2\\[2pt] \frac16(4) + 1\\[2pt] 1 \end{bmatrix}
= \begin{bmatrix} \frac23\\[2pt] \frac43\\[2pt] \frac53\\[2pt] 1 \end{bmatrix}
\]
\begin{align*}
x_1 &:\ 4 \div \tfrac23 = 6\\
s_2 &:\ 2 \div \tfrac43 = \mathbf{\tfrac32} \nota{la más chica}\\
s_3 &:\ 5 \div \tfrac53 = 3\\
s_4 &:\ 2 \div 1 = 2
\end{align*}
\sale{Sale $s_2$}: ahora es la materia prima M2 la que se agota.

\paso{5. Actualizo $\Binv$}
Sale la fila 2 con $\alpha_2 = \tfrac43$:
\[
\eta = \left(-\frac{2/3}{4/3},\ \frac{1}{4/3},\ -\frac{5/3}{4/3},\ -\frac{1}{4/3}\right)
= \left(-\tfrac12,\ \tfrac34,\ -\tfrac54,\ -\tfrac34\right)
\]
\[
E = \begin{bmatrix} 1 & -\frac12 & 0 & 0\\ 0 & \frac34 & 0 & 0\\ 0 & -\frac54 & 1 & 0\\ 0 & -\frac34 & 0 & 1\end{bmatrix},
\qquad
\Binv_{\text{nueva}} = E\,\Binv = \begin{bmatrix}
\frac14 & -\frac12 & 0 & 0\\ -\frac18 & \frac34 & 0 & 0\\ \frac38 & -\frac54 & 1 & 0\\ \frac18 & -\frac34 & 0 & 1
\end{bmatrix}.
\]
Por filas: la fila 2 $(-\tfrac16, 1, 0, 0)$ entre $\tfrac43$ da $(-\tfrac18, \tfrac34, 0, 0)$. A la fila 1 le resto
$\tfrac23$ de esa fila: $(\tfrac16, 0) - \tfrac23(-\tfrac18, \tfrac34) = (\tfrac14, -\tfrac12)$. Igual con las filas 3 y 4.

% ---------------------------------------------------------------------
\subsection*{Iteración 2 \quad (base $\{x_1, x_2, s_3, s_4\}$)}
\paso{1. Cuánto vale cada básica}
\begin{align*}
x_1 &= \tfrac14(24) - \tfrac12(6) = 6 - 3 = 3\\
x_2 &= -\tfrac18(24) + \tfrac34(6) = -3 + \tfrac92 = \tfrac32\\
s_3 &= \tfrac38(24) - \tfrac54(6) + 1 = 9 - \tfrac{15}{2} + 1 = \tfrac52\\
s_4 &= \tfrac18(24) - \tfrac34(6) + 2 = 3 - \tfrac92 + 2 = \tfrac12
\end{align*}
Con $c_B = (5,4,0,0)$: \quad $Z = 5(3) + 4\left(\tfrac32\right) = 15 + 6 = 21$.

\paso{2. Precios duales}
\begin{align*}
\pi_1 &= 5\cdot\tfrac14 + 4\cdot\left(-\tfrac18\right) = \tfrac54 - \tfrac12 = \tfrac34\\
\pi_2 &= 5\cdot\left(-\tfrac12\right) + 4\cdot\tfrac34 = -\tfrac52 + 3 = \tfrac12\\
\pi_3 &= 0, \qquad \pi_4 = 0
\end{align*}

\paso{3. ¿Alguien mejora $Z$?}
Las no básicas son $s_1$ y $s_2$:
\[
z_{s_1} - 0 = \pi_1 = \tfrac34 \ge 0, \qquad z_{s_2} - 0 = \pi_2 = \tfrac12 \ge 0.
\]
\textbf{Ninguno es negativo $\Rightarrow$ óptimo.}

\subsection*{Resultado}
\[
\res{x_1 = 3,\quad x_2 = \tfrac32,\quad Z^* = 21}
\]
\begin{itemize}[itemsep=2pt]
  \item Conviene producir \textbf{3 toneladas de exteriores y 1.5 de interiores} al día, ganando \textbf{\$21\,000}.
  \item M1 y M2 se usan completas ($s_1 = s_2 = 0$).
  \item $s_3 = \tfrac52$: la condición del mercado sobra mucho (interiores está 1.5 ton \emph{por debajo} de exteriores).
  \item $s_4 = \tfrac12$: se venden 1.5 ton de interiores, media tonelada menos que la demanda máxima.
  \item \textbf{Precios sombra} $(\tfrac34,\ \tfrac12,\ 0,\ 0)$: una tonelada extra de M1 subiría la ganancia en \$750,
        una de M2 en \$500. Las condiciones de mercado no están frenando nada (valen 0).
\end{itemize}

\subsection*{Interpretación gráfica}
\begin{center}
\begin{tikzpicture}[scale=1.05, >=Stealth]
  \fill[region!70] (0,0)--(4,0)--(3,1.5)--(2,2)--(1,2)--(0,1)--cycle;
  \draw[very thin, gray!25] (0,0) grid (6.5,4.5);
  \draw[->] (0,0)--(6.8,0) node[right] {$x_1$};
  \draw[->] (0,0)--(0,4.8) node[above] {$x_2$};
  \foreach \x in {1,...,6} \node[below, font=\scriptsize] at (\x,0) {\x};
  \foreach \y in {1,...,4} \node[left, font=\scriptsize] at (0,\y) {\y};
  \draw[r1] (1,4.5)--(4,0);
  \draw[r2] (0,3)--(6,0);
  \draw[r3] (0,1)--(3.5,4.5);
  \draw[r4] (0,2)--(6.5,2);
  \draw[zline] (0.6,4.5)--(4.2,0);
  \node[anchor=north east, draw=gray!50, fill=white, align=left, font=\scriptsize, inner sep=4pt] at (6.5,4.45) {
    \swatch{r1} R1: $6x_1 + 4x_2 \le 24$\\
    \swatch{r2} R2: $x_1 + 2x_2 \le 6$\\
    \swatch{r3} R3: $-x_1 + x_2 \le 1$\\
    \swatch{r4} R4: $x_2 \le 2$\\
    \swatch{zline} $Z = 21$};
  \draw[->, ultra thick, tinta!80] (0,0)--(3.85,0);
  \draw[->, ultra thick, tinta!80] (4,0)--(3.1,1.38);
  \foreach \p in {(0,0),(4,0),(2,2),(1,2),(0,1)} \fill[tinta] \p circle (2pt);
  \fill[resalta!90!black] (3,1.5) circle (3.5pt);
  \node[anchor=south west, fill=white, inner sep=1pt, font=\scriptsize\bfseries, xshift=3pt, yshift=3pt] at (0,0) {It.~0};
  \node[anchor=south west, fill=white, inner sep=1pt, font=\scriptsize\bfseries, xshift=3pt, yshift=3pt] at (4,0) {It.~1};
  \node[anchor=south west, fill=white, inner sep=1pt, font=\small\bfseries, xshift=5pt, yshift=5pt] at (3,1.5) {It.~2: $(3,\,1.5)$ óptimo};
\end{tikzpicture}
\end{center}
El simplex salió de $(0,0)$, avanzó por el eje hasta $(4,0)$ (se acabó M1) y subió a $(3,\,1.5)$ (se acabó M2).
Esquinas: $(0,0)\to 0$, \ $(4,0)\to 20$, \ $(3,1.5)\to \mathbf{21}$, \ $(2,2)\to 18$, \ $(1,2)\to 13$, \ $(0,1)\to 4$.

% =====================================================================
\newpage
\section{Terapia de radiación: minimizar con Gran M}
\textit{Fuente: Hillier y Lieberman, Introducción a la Investigación de Operaciones, cap.~3.}

\subsection*{Planteamiento del problema}
A una paciente con un tumor se le va a aplicar radiación con \textbf{dos haces}. Cada haz reparte su dosis entre
distintas zonas del cuerpo. Por cada kilorad que se dispara con cada haz, cada zona absorbe:

\begin{center}
\begin{tabular}{lccc}
\toprule
\textbf{Zona} & \textbf{Haz 1} & \textbf{Haz 2} & \textbf{Condición para esa zona} \\
\midrule
Tejido sano         & 0.4 & 0.5 & lo \textbf{mínimo} posible \\
Tejidos críticos    & 0.3 & 0.1 & como mucho 2.7 \\
Región del tumor    & 0.5 & 0.5 & exactamente 6 \\
Centro del tumor    & 0.6 & 0.4 & al menos 6 \\
\bottomrule
\end{tabular}
\end{center}

La pregunta: \textbf{¿qué dosis dar con cada haz para dañar lo menos posible el tejido sano?}

\subsection*{Formulación del modelo}
\begin{itemize}[itemsep=2pt]
  \item $x_1$ = dosis del haz 1 (kilorads), \quad $x_2$ = dosis del haz 2.
  \item \textbf{Objetivo:} lo que absorbe el tejido sano, $Z = 0.4x_1 + 0.5x_2$, y queremos que sea \textbf{mínimo}.
  \item \textbf{Tejidos críticos:} $0.3x_1 + 0.1x_2 \le 2.7$ \ (``como mucho'' $\to \le$).
  \item \textbf{Región del tumor:} $0.5x_1 + 0.5x_2 = 6$ \ (``exactamente'' $\to =$).
  \item \textbf{Centro del tumor:} $0.6x_1 + 0.4x_2 \ge 6$ \ (``al menos'' $\to \ge$).
\end{itemize}

\[
\begin{aligned}
\text{Min } Z &= 0.4x_1 + 0.5x_2\\
\text{s.a. } \quad 0.3x_1 + 0.1x_2 &\le 2.7\\
0.5x_1 + 0.5x_2 &= 6\\
0.6x_1 + 0.4x_2 &\ge 6\\
x_1, x_2 &\ge 0
\end{aligned}
\]
Para trabajar exacto uso fracciones: $0.4 = \tfrac25$, $0.5 = \tfrac12$, $0.3 = \tfrac{3}{10}$, $0.1 = \tfrac1{10}$,
$0.6 = \tfrac35$, $2.7 = \tfrac{27}{10}$.

\subsection*{Forma estándar}
\begin{enumerate}[itemsep=3pt]
  \item \textbf{Min a Max:} minimizar $Z$ es lo mismo que maximizar $-Z$. Así que trabajo con
        $\text{Max } (-Z) = -\tfrac25 x_1 - \tfrac12 x_2$ y al final le cambio el signo.
  \item \textbf{Restricción $\le$:} le sumo la holgura $s_1$, igual que antes.
  \item \textbf{Restricción $=$:} ya es igualdad, pero \textbf{no tiene ninguna variable que pueda arrancar en la base}.
        Le invento una \textbf{artificial} $a_2$.
  \item \textbf{Restricción $\ge$:} le resto un \textbf{exceso} $e_3$ (lo que me paso de 6). Pero $e_3$ tiene coeficiente $-1$
        y no sirve para la base inicial (valdría $-6$), así que también le sumo una artificial $a_3$.
\end{enumerate}
\begin{align*}
\tfrac{3}{10}x_1 + \tfrac1{10}x_2 + s_1 &= \tfrac{27}{10}\\
\tfrac12 x_1 + \tfrac12 x_2 + a_2 &= 6\\
\tfrac35 x_1 + \tfrac25 x_2 - e_3 + a_3 &= 6
\end{align*}

\begin{idea}[Método de la Gran M]
Las variables artificiales no forman parte del problema real: solo existen para tener una base al inicio. En la solución final
\textbf{tienen que valer 0}. Para obligar al método a sacarlas, las castigo en la función objetivo con $-M$,
donde $M$ es un número gigantesco:
\[
\text{Max } (-Z) = -\tfrac25 x_1 - \tfrac12 x_2 + 0s_1 - M a_2 + 0e_3 - M a_3.
\]
Regla para comparar: \textbf{manda el término con $M$}. Solo si dos cosas tienen la misma cantidad de $M$,
se desempata con el número suelto.
\end{idea}

Orden de columnas $(x_1, x_2, s_1, a_2, e_3, a_3)$, con $b = \left(\tfrac{27}{10},\ 6,\ 6\right)$ y
$c = \left(-\tfrac25,\ -\tfrac12,\ 0,\ -M,\ 0,\ -M\right)$. Base inicial $\{s_1, a_2, a_3\}$, $\Binv = I$.

% ---------------------------------------------------------------------
\subsection*{Iteración 0 \quad (base $\{s_1, a_2, a_3\}$)}
\paso{1 y 2. Valores y precios duales}
\[
x_B = \left(\tfrac{27}{10},\ 6,\ 6\right), \qquad c_B = (0,\ -M,\ -M), \qquad \pi = c_B\,I = (0,\ -M,\ -M)
\]
\[
(-Z) = 0\cdot\tfrac{27}{10} - 6M - 6M = -12M \nota{muy malo: las artificiales valen 6 cada una}
\]

\paso{3. ¿Alguien mejora?}
Las no básicas son $x_1$, $x_2$ y $e_3$:
\begin{align*}
z_1 - c_1 &= \left[0\cdot\tfrac{3}{10} - M\cdot\tfrac12 - M\cdot\tfrac35\right] + \tfrac25 = \tfrac25 - \tfrac{11}{10}M \nota{el más negativo}\\
z_2 - c_2 &= \left[-M\cdot\tfrac12 - M\cdot\tfrac25\right] + \tfrac12 = \tfrac12 - \tfrac{9}{10}M\\
z_{e_3} - c_{e_3} &= -M\cdot(-1) - 0 = M
\end{align*}
Comparo solo la $M$: $-\tfrac{11}{10}M$ es más negativo que $-\tfrac{9}{10}M$. \entra{Entra $x_1$}.

\paso{4. ¿Quién sale?}
\[
\alpha = a_1 = \left(\tfrac{3}{10},\ \tfrac12,\ \tfrac35\right)
\]
\begin{align*}
s_1 &:\ \tfrac{27}{10} \div \tfrac{3}{10} = \mathbf{9}\\
a_2 &:\ 6 \div \tfrac12 = 12\\
a_3 &:\ 6 \div \tfrac35 = 10
\end{align*}
\sale{Sale $s_1$}: los tejidos críticos llegan a su límite con $x_1 = 9$.

\paso{5. Actualizo $\Binv$}
Sale la fila 1 con $\alpha_1 = \tfrac{3}{10}$:
$\eta = \left(\tfrac{10}{3},\ -\tfrac{1/2}{3/10},\ -\tfrac{3/5}{3/10}\right) = \left(\tfrac{10}3,\ -\tfrac53,\ -2\right)$.
\[
\Binv_{\text{nueva}} = E = \begin{bmatrix} \frac{10}{3} & 0 & 0\\ -\frac53 & 1 & 0\\ -2 & 0 & 1\end{bmatrix}
\]

% ---------------------------------------------------------------------
\subsection*{Iteración 1 \quad (base $\{x_1, a_2, a_3\}$)}
\paso{1. Valores}
\begin{align*}
x_1 &= \tfrac{10}{3}\cdot\tfrac{27}{10} = 9\\
a_2 &= -\tfrac53\cdot\tfrac{27}{10} + 6 = -\tfrac92 + 6 = \tfrac32 \nota{la artificial bajó de 6 a 1.5}\\
a_3 &= -2\cdot\tfrac{27}{10} + 6 = -\tfrac{27}{5} + 6 = \tfrac35
\end{align*}
$(-Z) = -\tfrac25(9) - M\cdot\tfrac32 - M\cdot\tfrac35 = -\tfrac{18}{5} - \tfrac{21}{10}M$ \quad (mejoró: antes era $-12M$).

\paso{2. Precios duales} \ con $c_B = \left(-\tfrac25,\ -M,\ -M\right)$:
\begin{align*}
\pi_1 &= -\tfrac25\cdot\tfrac{10}{3} - M\cdot\left(-\tfrac53\right) - M\cdot(-2) = -\tfrac43 + \tfrac{11}{3}M\\
\pi_2 &= -M, \qquad \pi_3 = -M
\end{align*}

\paso{3. ¿Alguien mejora?}
\begin{align*}
z_2 - c_2 &= \pi_1\cdot\tfrac1{10} + \pi_2\cdot\tfrac12 + \pi_3\cdot\tfrac25 + \tfrac12 = \tfrac{11}{30} - \tfrac{8}{15}M \nota{negativo}\\
z_{s_1} - c_{s_1} &= \pi_1 = -\tfrac43 + \tfrac{11}{3}M \nota{la $M$ positiva manda: positivo}\\
z_{e_3} - c_{e_3} &= -\pi_3 = M
\end{align*}
\entra{Entra $x_2$}.

\paso{4. ¿Quién sale?}
\[
\alpha = \Binv a_2 = \begin{bmatrix} \frac{10}{3}\cdot\frac1{10}\\[2pt] -\frac53\cdot\frac1{10} + \frac12\\[2pt] -2\cdot\frac1{10} + \frac25\end{bmatrix}
= \begin{bmatrix} \frac13\\[2pt] \frac13\\[2pt] \frac15\end{bmatrix}
\]
\begin{align*}
x_1 &:\ 9 \div \tfrac13 = 27\\
a_2 &:\ \tfrac32 \div \tfrac13 = \tfrac92\\
a_3 &:\ \tfrac35 \div \tfrac15 = \mathbf{3} \nota{la más chica}
\end{align*}
\sale{Sale $a_3$}: ¡ya se fue la primera artificial!

\paso{5. Actualizo $\Binv$}
Sale la fila 3 con $\alpha_3 = \tfrac15$: \ $\eta = \left(-\tfrac53,\ -\tfrac53,\ 5\right)$.
\[
E = \begin{bmatrix} 1 & 0 & -\frac53\\ 0 & 1 & -\frac53\\ 0 & 0 & 5\end{bmatrix},
\qquad
\Binv_{\text{nueva}} = \begin{bmatrix} \frac{20}{3} & 0 & -\frac53\\ \frac53 & 1 & -\frac53\\ -10 & 0 & 5\end{bmatrix}
\]
Por filas: la fila 3 $(-2, 0, 1)$ por 5 da $(-10, 0, 5)$. A la fila 1 le resto $\tfrac13$ de eso:
$(\tfrac{10}3, 0, 0) - \tfrac13(-10, 0, 5) = (\tfrac{20}{3}, 0, -\tfrac53)$. Lo mismo con la fila 2.

% ---------------------------------------------------------------------
\subsection*{Iteración 2 \quad (base $\{x_1, a_2, x_2\}$)}
\paso{1. Valores}
\begin{align*}
x_1 &= \tfrac{20}{3}\cdot\tfrac{27}{10} - \tfrac53\cdot 6 = 18 - 10 = 8\\
a_2 &= \tfrac53\cdot\tfrac{27}{10} + 6 - \tfrac53\cdot 6 = \tfrac92 + 6 - 10 = \tfrac12\\
x_2 &= -10\cdot\tfrac{27}{10} + 5\cdot 6 = -27 + 30 = 3
\end{align*}
$(-Z) = -\tfrac25(8) - \tfrac12 M - \tfrac12(3) = -\tfrac{47}{10} - \tfrac12 M$.

\paso{2. Precios duales} \ con $c_B = \left(-\tfrac25,\ -M,\ -\tfrac12\right)$:
\begin{align*}
\pi_1 &= -\tfrac25\cdot\tfrac{20}3 - M\cdot\tfrac53 - \tfrac12\cdot(-10) = \tfrac73 - \tfrac53 M\\
\pi_2 &= -M\\
\pi_3 &= -\tfrac25\cdot\left(-\tfrac53\right) - M\cdot\left(-\tfrac53\right) - \tfrac12\cdot 5 = -\tfrac{11}{6} + \tfrac53 M
\end{align*}

\paso{3. ¿Alguien mejora?}
\begin{align*}
z_{s_1} - c_{s_1} &= \pi_1 = \tfrac73 - \tfrac53 M \nota{negativo}\\
z_{e_3} - c_{e_3} &= -\pi_3 = \tfrac{11}{6} - \tfrac53 M \nota{negativo}\\
z_{a_3} - c_{a_3} &= \pi_3 + M = -\tfrac{11}6 + \tfrac83 M \nota{positivo}
\end{align*}
Aquí hay un \textbf{empate en la $M$}: los dos tienen $-\tfrac53 M$. Desempata el número suelto:
$\tfrac{11}{6} < \tfrac73$, así que \entra{entra $e_3$}.

\paso{4. ¿Quién sale?}
La columna de $e_3$ es $(0,\ 0,\ -1)$, así que $\alpha$ es la columna 3 de $\Binv$ cambiada de signo:
\[
\alpha = \left(\tfrac53,\ \tfrac53,\ -5\right)
\]
\begin{align*}
x_1 &:\ 8 \div \tfrac53 = \tfrac{24}{5}\\
a_2 &:\ \tfrac12 \div \tfrac53 = \mathbf{\tfrac{3}{10}} \nota{la más chica}\\
x_2 &:\ \alpha = -5 \ \text{no cuenta}
\end{align*}
\sale{Sale $a_2$}: ¡se va la última artificial! Desde aquí ya estamos dentro de la región factible.

\paso{5. Actualizo $\Binv$}
Sale la fila 2 con $\alpha_2 = \tfrac53$: \ $\eta = \left(-1,\ \tfrac35,\ 3\right)$.
\[
E = \begin{bmatrix} 1 & -1 & 0\\ 0 & \frac35 & 0\\ 0 & 3 & 1\end{bmatrix},
\qquad
\Binv_{\text{nueva}} = \begin{bmatrix} 5 & -1 & 0\\ 1 & \frac35 & -1\\ -5 & 3 & 0\end{bmatrix}
\]

% ---------------------------------------------------------------------
\subsection*{Iteración 3 \quad (base $\{x_1, e_3, x_2\}$)}
\paso{1. Valores}
\begin{align*}
x_1 &= 5\cdot\tfrac{27}{10} - 6 = \tfrac{27}{2} - 6 = \tfrac{15}{2}\\
e_3 &= \tfrac{27}{10} + \tfrac35\cdot 6 - 6 = \tfrac{27}{10} + \tfrac{36}{10} - \tfrac{60}{10} = \tfrac{3}{10}\\
x_2 &= -5\cdot\tfrac{27}{10} + 3\cdot 6 = -\tfrac{27}{2} + 18 = \tfrac92
\end{align*}
Con $c_B = \left(-\tfrac25,\ 0,\ -\tfrac12\right)$: \quad
$(-Z) = -\tfrac25\cdot\tfrac{15}2 - \tfrac12\cdot\tfrac92 = -3 - \tfrac94 = -\tfrac{21}{4}$. \ ¡Ya no hay $M$!

\paso{2. Precios duales}
\begin{align*}
\pi_1 &= -\tfrac25(5) + 0(1) - \tfrac12(-5) = -2 + \tfrac52 = \tfrac12\\
\pi_2 &= -\tfrac25(-1) + 0\left(\tfrac35\right) - \tfrac12(3) = \tfrac25 - \tfrac32 = -\tfrac{11}{10}\\
\pi_3 &= -\tfrac25(0) + 0(-1) - \tfrac12(0) = 0
\end{align*}

\paso{3. ¿Alguien mejora?}
\begin{align*}
z_{s_1} - c_{s_1} &= \pi_1 = \tfrac12 \ge 0\\
z_{a_2} - c_{a_2} &= \pi_2 + M = -\tfrac{11}{10} + M \ge 0 \nota{$M$ gigante}\\
z_{a_3} - c_{a_3} &= \pi_3 + M = M \ge 0
\end{align*}
\textbf{Ninguno es negativo $\Rightarrow$ óptimo.}

\subsection*{Resultado}
Como trabajamos con $\text{Max}(-Z) = -\tfrac{21}{4}$, el mínimo es $Z^* = \tfrac{21}{4}$:
\[
\res{x_1 = \tfrac{15}{2} = 7.5,\quad x_2 = \tfrac92 = 4.5,\quad Z^* = \tfrac{21}{4} = 5.25}
\]
\begin{itemize}[itemsep=2pt]
  \item Hay que dar \textbf{7.5 kilorads con el haz 1 y 4.5 con el haz 2}. El tejido sano absorbe \textbf{5.25}, lo mínimo posible.
  \item $s_1 = 0$: los tejidos críticos quedan justo en su límite: $0.3(7.5) + 0.1(4.5) = 2.7$.
  \item $e_3 = \tfrac{3}{10}$: el centro del tumor recibe $0.6(7.5) + 0.4(4.5) = 6.3$, o sea 0.3 más del mínimo.
  \item Las artificiales terminaron en 0, como debía ser: la solución sí es factible.
  \item \textbf{Precios sombra} del problema original: $-\tfrac12$ para los tejidos críticos (si se permitiera 1 kilorad más ahí,
        el daño al tejido sano \emph{bajaría} 0.5), $\tfrac{11}{10}$ para la región del tumor (si se exigiera 1 más,
        el daño \emph{subiría} 1.1) y 0 para el centro del tumor (esa condición no está frenando nada).
\end{itemize}

\subsection*{Interpretación gráfica}
\begin{center}
\begin{tikzpicture}[scale=0.52, >=Stealth]
  \draw[very thin, gray!25] (0,0) grid (14,16);
  \draw[->] (0,0)--(14.6,0) node[right] {$x_1$};
  \draw[->] (0,0)--(0,16.6) node[above] {$x_2$};
  \foreach \x in {2,4,...,14} \node[below, font=\scriptsize] at (\x,0) {\x};
  \foreach \y in {2,4,...,16} \node[left, font=\scriptsize] at (0,\y) {\y};
  \draw[r1] (9,0)--(3.667,16);
  \draw[r2] (0,12)--(12,0);
  \draw[r3] (0,15)--(10,0);
  \draw[zline] (0,10.5)--(13.125,0);
  \node[anchor=north east, draw=gray!50, fill=white, align=left, font=\scriptsize, inner sep=4pt] at (14,15.8) {
    \swatch{r1} R1: $0.3x_1 + 0.1x_2 \le 2.7$\\
    \swatch{r2} R2: $0.5x_1 + 0.5x_2 = 6$\\
    \swatch{r3} R3: $0.6x_1 + 0.4x_2 \ge 6$\\
    \swatch{zline} $Z = 5.25$};
  \draw[line width=4pt, region!90!black, line cap=round] (6,6)--(7.5,4.5);
  \draw[->, ultra thick, tinta!80] (0,0)--(8.8,0);
  \draw[->, ultra thick, tinta!80] (9,0)--(8.07,2.8);
  \draw[->, ultra thick, tinta!80] (8,3)--(7.57,4.35);
  \fill[tinta] (6,6) circle (3pt);
  \foreach \p in {(0,0),(9,0),(8,3)} \draw[fill=white, draw=tinta] \p circle (5pt);
  \fill[resalta!90!black] (7.5,4.5) circle (8pt);
  \node[anchor=south west, fill=white, inner sep=1pt, font=\scriptsize\bfseries, xshift=3pt, yshift=3pt] at (0,0) {It.~0};
  \node[anchor=south west, fill=white, inner sep=1pt, font=\scriptsize\bfseries, xshift=3pt, yshift=3pt] at (9,0) {It.~1};
  \node[anchor=south west, fill=white, inner sep=1pt, font=\scriptsize\bfseries, xshift=3pt, yshift=3pt] at (8,3) {It.~2};
  \node[anchor=south west, fill=white, inner sep=1pt, font=\small\bfseries, xshift=5pt, yshift=5pt] at (7.5,4.5) {It.~3: $(7.5,\,4.5)$ óptimo};
\end{tikzpicture}
\end{center}
Como hay una igualdad, la región factible no es un área sino \textbf{el pedacito de recta} entre $(6,6)$ y $(7.5,\,4.5)$
(la línea gruesa). Fíjate en los círculos blancos: con Gran M el simplex pasa por puntos que \textbf{no son factibles}
($(0,0)$, $(9,0)$, $(8,3)$), porque las artificiales todavía valen algo. La $M$ lo va empujando hasta que entra a la región
y termina en $(7.5,\,4.5)$. Evaluando los dos extremos: $(6,6)\to 5.4$ y $(7.5,4.5)\to \mathbf{5.25}$. Coincide.

% =====================================================================
\newpage
\section*{Resumen de los tres ejercicios}
\begin{center}
\begin{tabular}{lccccc}
\toprule
\textbf{Ejercicio} & \textbf{Tipo} & \textbf{Técnica} & \textbf{Iteraciones} & \textbf{Solución} & \textbf{Óptimo} \\
\midrule
Wyndor Glass & Max, todo $\le$ & holguras & 2 & $(2,\ 6)$ & $Z = 36$ \\
Reddy Mikks & Max, todo $\le$ & holguras & 2 & $(3,\ 1.5)$ & $Z = 21$ \\
Radiación & Min, $\le$ $=$ $\ge$ & Gran M & 3 & $(7.5,\ 4.5)$ & $Z = 5.25$ \\
\bottomrule
\end{tabular}
\end{center}

\begin{idea}[Conclusiones]
\begin{itemize}[itemsep=2pt, leftmargin=*]
  \item En el Simplex Revisado todo se obtiene de $\Binv$: con ella se calculan $x_B$, $\pi$ y $\alpha$ sin actualizar la tabla completa.
  \item \entra{Entra} la no básica con $z_j - c_j$ más negativo; \sale{sale} la básica con la menor razón $x_B/\alpha$ (solo $\alpha > 0$).
  \item $\Binv$ se actualiza con $E$, que es hacer Gauss sobre la fila del pivote.
  \item Cuando hay $=$ o $\ge$ aparecen artificiales; la Gran M las castiga para que salgan de la base.
  \item Los $\pi$ del final son los \textbf{precios sombra}: cuánto cambia el óptimo si tuviera una unidad más de cada recurso.
\end{itemize}
\end{idea}

\textbf{Verificación:} los tres ejercicios se comprobaron con la implementación del método (\emph{Simplex Arcade},
modo 1P, Gran M), que reproduce las mismas iteraciones, y con la librería javascript-lp-solver, que confirma los tres óptimos.

\section*{Referencias}
\addcontentsline{toc}{part}{Referencias}
\begin{itemize}[leftmargin=*, itemsep=3pt]
  \item Bazaraa, M. S., Jarvis, J. J. y Sherali, H. D. \textit{Programación lineal y flujo en redes}. Limusa.
  \item Dantzig, G. B. y Orchard-Hays, W. (1954). The product form for the inverse in the simplex method. \textit{Mathematical Tables and Other Aids to Computation}, 8(46).
  \item Hillier, F. S. y Lieberman, G. J. \textit{Introducción a la Investigación de Operaciones}. McGraw-Hill.
  \item Taha, H. A. \textit{Investigación de Operaciones}. Pearson.
\end{itemize}

\end{document}
